\section{Cross-phase acceptance matrix}

The following matrix prevents an implementation or provider from being declared complete on the basis of a successful build alone.
\begin{longtable}{@{}L{0.22\textwidth}L{0.16\textwidth}L{0.16\textwidth}p{0.35\textwidth}@{}}
\toprule
\textbf{Evidence} & \textbf{Core phase} & \textbf{Package phase} & \textbf{Required proof} \\
\midrule
API/compile compatibility & 1--3 & Every package & Public contracts compile; one operation type can have multiple typed implementation classes; no backend-native types leak. \\
Semantic operation tests & 1, 3 & Every operation implementation & Inputs, outputs, numerical rules, layer semantics, and edge cases satisfy the operation specification. \\
Generic lifecycle tests & 2, 4, 6 & Every provider & Resource ownership, completion, synchronization, failure, cancellation, reservation, and release are correct. \\
Specification tests & 1, 3, 4 & Every operation implementation & Declared support matches actual tested behavior; unsupported modes and version ranges are rejected. \\
Policy-construction tests & 1, 3 & Core protocol/runtime & Contradictory requirements are rejected at construction/composition with provenance-rich diagnostics. \\
Resolver tests & 4--5 & Every implementation candidate & Correct eligibility and ranking using operation, implementation spec, environment state, and policy; no invalid candidate is selected. \\
Dynamic-state tests & 2, 4--6 & Every observable provider & Busy/unavailable devices, memory pressure, stale telemetry, and reservation races are handled safely. \\
Image transaction tests & 1, 3, 6 & Every operation implementation & New-image and replacement semantics, identity, revisions, and failure atomicity are preserved. \\
Real-device tests & N/A & GPU providers & Actual device execution, resource sharing, precision, synchronization, and errors are verified. \\
Performance tests & 5--8 & Optimized packages & End-to-end benefit is measured with queueing, transfer, allocation, warm-up, and synchronization included. \\
\bottomrule
\end{longtable}

\section{Initial operation scope and ordering}

The first supported operation set should be intentionally small. The suggested order is:
\begin{enumerate}
  \item \textbf{Descriptor-preserving conversion:} explicit pixel/component format conversion, channel mapping, and numeric conversion. Conversions state loss and range behavior.
  \item \textbf{Resize:} specifies interpolation, coordinate convention, border handling, and label-aware restrictions.
  \item \textbf{Gaussian filtering:} specifies kernel generation, sigma/radius relationship, border extension, computation/accumulation/output precision, and channel/alpha behavior.
  \item \textbf{Color correction:} specifies color space, transfer function, numeric range, and whether input is color data rather than scientific intensity.
  \item \textbf{Compositing/blending:} specifies blend modes, alpha convention, working space, opacity, and layer ordering.
\end{enumerate}

The first end-to-end resolver demonstration should use one operation with multiple implementation classes. Gaussian blur is a suitable vertical slice because its border behavior, kernel parameters, precision, and numerical tolerances expose meaningful implementation differences.

Annotation is initially a typed layer and rendering concern, not an excuse to treat all annotations as raster pixels. Tracking is a stateful/temporal operation family and should follow explicit sequence inputs, state ownership, reset behavior, and temporal error semantics.

\section{Release strategy and support claims}

A release is a tested matrix, not a list of compiled adapters. For every released version, publish:
\begin{itemize}
  \item supported OS, architecture, compiler/toolchain, and runtime versions;
  \item provider/device combinations actually tested;
  \item operation/implementation combinations that pass conformance tests;
  \item supported precision, format, layout, parameter, and semantic constraints;
  \item known unsupported cases and fallback behavior;
  \item which execution-state metrics are available, estimated, or unknown;
  \item build/test commands and required SDKs/drivers;
  \item benchmark methodology and measured workloads, without generalizing one device's result to all devices.
\end{itemize}

A provider can be experimental while its protocol implementation is being validated. It must not be labeled stable because a basic sample ran once. An algorithm adapter may support only a subset of operation parameters. Implementation specifications, resolver diagnostics, and the published support matrix must reflect that subset.

\section{Risks, stop conditions, and scope control}

\begin{longtable}{@{}L{0.22\textwidth}p{0.69\textwidth}@{}}
\toprule
\textbf{Risk / stop condition} & \textbf{Required response} \\
\midrule
The generic execution protocol begins accumulating image terms & Stop and move the image-specific rule into \IPP{}. Keep generic resource, provider, capability, state-observation, submit, and synchronization concepts only. \\
The image protocol exposes native Metal/CUDA/OpenCV types & Stop and redesign around semantic descriptors plus opaque generic resources. \\
An operation specification is too vague to compare implementations & Do not register multiple implementations. Specify semantics and reference tests first. \\
An implementation spec overstates actual support & Reject registration or qualification until declaration and conformance evidence agree. \\
The resolver embeds one-off provider/platform branches & Replace them with declarative implementation requirements and generic environment capabilities. \\
Unknown device telemetry is treated as idle & Represent unknown explicitly and use other measurements without claiming certainty. \\
A provider cannot satisfy required numerical semantics & Reject it for that operation/mode; do not silently lower precision. \\
A provider requires hidden copies or materialization & Make the transition explicit in qualification/planning and measure its cost. \\
A required provider is unavailable & Report an unsatisfied requirement; do not silently fall back. \\
A preference conflict is silently resolved by declaration order & Reject ambiguous same-scope objective configuration or require an explicit priority/weight model. \\
Backend-specific behavior leaks into callers & Add declarative capability/qualification information or a legitimate operation parameter; do not scatter platform checks through application code. \\
Optimization has no measured end-to-end benefit & Defer it. Keep the simpler execution path. \\
A provider lacks real hardware CI & Keep it experimental; do not claim production support. \\
Scope expands into training, UI, or unrelated workflow orchestration & Keep those concerns outside the core library unless a concrete integration contract is approved. \\
\bottomrule
\end{longtable}

\section{Phase completion record}

Each phase should be closed with a short record containing:
\begin{enumerate}
  \item phase identifier and scope completed;
  \item public/API changes and decisions made;
  \item commit references and package versions;
  \item build and test commands with platform/hardware details;
  \item acceptance evidence and known failures;
  \item deviations from this plan and their rationale;
  \item explicit approval to start the next phase.
\end{enumerate}

\section{Final implementation principle}

Build the operation/domain contracts first; implement the generic execution dependency without allowing it to redefine the domain; then build the registry, environment discovery, common resolver, and execution coordinator. Keep concrete providers and operation algorithms in external packages selected by the application. Establish conformance on a small reference workload, and optimize only measured end-to-end behavior.

The objective is not to support the largest number of APIs. It is to preserve one image-processing meaning across independently supplied implementations that the resolver can qualify, rank, and execute safely.